\subsection{Numerical Methods}
\label{subsection:numerical-methods}

    In our project, we implemented three numerical methods for solving the H-H equations: Euler's method, improved Euler's method, and the fourth-order Runge-Kutta method.  
    We outline the mathematical formulae for each method below.

    \subsubsection{Euler's Method}
    \label{subsubsec:euler-method}

        Euler's method is a straightforward numerical technique for solving differential equations by constructing approximations using tangent lines.  
        Given the initial value problem:

        \[
        \frac{dy}{dt} = f(t, y), \quad y(t_0) = y_0,
        \]

        we start at $(t_0, y_0)$ and use the slope to form the tangent line.  
        The iterative formula is:

        \[
        y_{n+1} = y_n + f(t_n, y_n) \cdot \Delta t,
        \]

        where $\Delta t = t_{n+1} - t_n$ is the step size.  
        As $\Delta t \to 0$, this piecewise linear approximation converges to the true solution $\phi(t)$.

    \subsubsection{Improved Euler's Method}
    \label{subsubsec:improved-euler}

        The improved Euler's method (Heun's method) is a predictor-corrector approach that achieves better accuracy by averaging slopes.  
        Given the same initial value problem, the method proceeds in two steps:

        \[
        \begin{aligned}
        y_{n+1}^* &= y_n + f(t_n, y_n) \Delta t\\[6pt]
        y_{n+1}   &= y_n + \frac{f(t_n, y_n) + f(t_{n+1}, y_{n+1}^*)}{2} \Delta t
        \end{aligned}
        \]

        The predictor uses standard Euler, then the corrector averages the slopes at the beginning and end of the interval for better accuracy.  
        This method provides second-order accuracy compared to Euler's first-order approximation.

    \subsubsection{Runge-Kutta Method}
    \label{subsubsec:runge-kutta}

        The fourth-order Runge-Kutta method provides excellent accuracy with reasonable computational cost.  
        It requires four slope evaluations per step:

        \[
        \begin{aligned}
        k_1 &= f(t_n, y_n) \\[3pt]
        k_2 &= f\bigl(t_n + \tfrac{h}{2},\; y_n + \tfrac{h}{2}k_1\bigr) \\[3pt]
        k_3 &= f\bigl(t_n + \tfrac{h}{2},\; y_n + \tfrac{h}{2}k_2\bigr) \\[3pt]
        k_4 &= f\bigl(t_n + h,\; y_n + h k_3\bigr)
        \end{aligned}
        \]

        where $h = \Delta t$.  
        These slope estimates are combined as:

        \[
        y_{n+1} = y_n + \frac{h}{6}(k_1 + 2k_2 + 2k_3 + k_4).
        \]

        Despite more function evaluations per step, its higher accuracy often allows larger step sizes.

\subsection{Noise Generation}
\label{subsection:noise-generation}

    To produce more realistic synthetic data, three key noise types were added to voltage signals: (1) Gaussian noise, (2) amplitude modulation, and (3) ocular artifacts.  
    Gaussian noise was implemented by adding random samples from a zero-mean normal distribution:

    \[
    \mathcal{N}(0, \sigma^2)
    \]

    with specified standard deviation $\sigma$.  
    Amplitude modulation applied slow amplitude variations using a tapered cosine window, simulating instrumental drift or movement artifacts.  
    Ocular artifacts introduced sharp, triangular-shaped perturbations to mimic abrupt physiological interference.  
    These methods were adapted from EEG noise models but scaled to action potential amplitudes \citep{SEREEGA2018}.

\subsection{Finite Impulse Response (FIR) Filter}
\label{subsection:fir-filter}

    A Finite Impulse Response (FIR) filter computes the output $y[n]$ through time-domain convolution of the input signal $x[n]$ with a finite set of filter coefficients $h[k]$:

    \[
    y[n] = \sum_{k=0}^{M-1} h[k] \cdot x[n-k]
    \]

    where $M$ is the filter order.  
    We implemented 58th-order FIR filter that uses a circular buffer to efficiently compute the convolution in the time domain.  
    Rather than shifting the entire input array, the algorithm maintains a buffer of length $M$ and updates a circular pointer at each time step.  
    The output at each step is computed as the dot product of the buffer and the coefficient vector.  

\subsection{Spike Detection}
\label{subsection:spike-detection}

    Spike detection was implemented using a percentile-based thresholding method with refractory period enforcement.  
    The threshold $\theta$ was calculated as the 95th percentile of the voltage signal:

    \[
    \theta = P_{95}(V)
    \]

    where $V$ represents the voltage time series.  
    To avoid noise-driven false positives, baseline noise was estimated from the lower 80\% of voltage values, with standard deviation $\sigma_{\text{noise}}$.  
    A spike was detected when the voltage crossed $\theta$ from below.  
    After each detection, a refractory period of $t_{\text{refract}}$ milliseconds was enforced.  
    Given a sampling rate $f_s$, the refractory period in samples is:

    \[
    n_{\text{refract}} = f_s \cdot \frac{t_{\text{refract}}}{1000}.
    \]

    The algorithm iterates through the voltage trace, identifying threshold crossings that satisfy the refractory constraint, and returns spike times in milliseconds.